\section{Methodology}

Our experimental design systematically verifies the causal mechanism linking CoT prompting to improved confidence calibration. Rather than simply measuring end-to-end ECE improvement, we decompose the hypothesized mechanism into four testable steps and validate each independently.

\subsection{Mechanism Hypothesis}

We hypothesize that CoT+confidence improves calibration through a ``self-reading'' mechanism:

\begin{enumerate}
    \item CoT prompting forces the model to articulate reasoning before answering
    \item This reasoning reveals epistemic uncertainty through hedging markers
    \item Hedging markers appear in-context before confidence generation (due to autoregressive generation)
    \item The model incorporates these uncertainty signals into its confidence estimate
\end{enumerate}

This mechanism predicts a negative correlation between hedging marker count and verbalized confidence: more hedging should lead to lower confidence.

\subsection{Sub-Hypothesis Decomposition}

We decompose the main hypothesis into four mechanism sub-hypotheses, each testing a specific component:

\textbf{H-M1 (Mechanism Step 1):} CoT prompting produces multi-step reasoning chains ($>$90\% reasoning rate, mean steps $>$2.0).

\textbf{H-M2 (Mechanism Step 2):} CoT outputs contain epistemic hedging markers ($>$30\% presence rate).

\textbf{H-M3 (Mechanism Step 3):} Hedging markers appear before confidence verbalization in the output sequence ($>$99\% structural compliance).

\textbf{H-M4 (Mechanism Step 4):} Hedging marker count negatively correlates with verbalized confidence (Spearman $r < -0.2$, $p < 0.05$).

\subsection{Experimental Design}

\textbf{Dataset:} We use TruthfulQA \citep{lin2022truthfulqa}, an adversarial benchmark testing model tendency toward common misconceptions. The generation split contains 817 items, providing sufficient statistical power for correlation analysis.

\textbf{Model:} GPT-3.5-turbo via OpenAI API, with temperature=0 for deterministic outputs. This represents a widely-deployed instruction-tuned model capable of following structured prompting.

\textbf{Prompting Conditions:} Our primary evaluation uses the CoT+Confidence condition: CoT reasoning + answer + confidence. This is compared against a baseline condition (direct answer only) for reasoning rate verification.

\textit{Note:} Our experimental design originally included a token-padding control condition (random filler tokens + answer + confidence) to rule out token-count artifacts. This condition was not executed due to resource constraints and remains future work. The current study validates the mechanism but does not empirically rule out token-count confounds.

\subsection{Measurement Protocols}

\textbf{Hedging Marker Detection:} We identify 17 epistemic hedging markers based on linguistic literature \citep{hyland1998hedging}: ``may,'' ``could,'' ``might,'' ``possibly,'' ``perhaps,'' ``likely,'' ``unlikely,'' ``but,'' ``however,'' ``although,'' ``alternatively,'' ``uncertain,'' ``not sure,'' ``seems,'' ``appears,'' ``probably,'' ``suggest.''

\textbf{Confidence Extraction:} We parse verbalized confidence from outputs using regex matching for patterns like ``Confidence: X\%'' with fallback to percentage mentions.

\textbf{Positional Analysis:} We verify structural ordering by identifying CoT, answer, and confidence segments and confirming markers appear before confidence statements.

\textbf{Correlation Analysis:} We compute Spearman rank correlation between hedging count and confidence, with significance testing via permutation.
